\documentclass{amsart}
% For dealing with references we use the comment environment
\usepackage{verbatim}
\newenvironment{reference}{\comment}{\endcomment}
%\newenvironment{reference}{}{}
\newenvironment{slogan}{\comment}{\endcomment}
\newenvironment{history}{\comment}{\endcomment}

% For commutative diagrams we use Xy-pic
\usepackage[all]{xy}

% We use 2cell for 2-commutative diagrams.
\xyoption{2cell}
\UseAllTwocells

% We use multicol for the list of chapters between chapters
\usepackage{multicol}

% This is generally recommended for better output
\usepackage{lmodern}
\usepackage[T1]{fontenc}

% For cross-file-references

% Package for hypertext links:
\usepackage{hyperref}

% For any local file, say "hello.tex" you want to link to please

% Theorem environments.
%
\theoremstyle{plain}
\newtheorem{theorem}[subsection]{Theorem}
\newtheorem{proposition}[subsection]{Proposition}
\newtheorem{lemma}[subsection]{Lemma}

\theoremstyle{definition}
\newtheorem{definition}[subsection]{Definition}
\newtheorem{example}[subsection]{Example}
\newtheorem{exercise}[subsection]{Exercise}
\newtheorem{situation}[subsection]{Situation}

\theoremstyle{remark}
\newtheorem{remark}[subsection]{Remark}
\newtheorem{remarks}[subsection]{Remarks}

\numberwithin{equation}{subsection}

% Macros
%
\def\lim{\mathop{\mathrm{lim}}\nolimits}
\def\colim{\mathop{\mathrm{colim}}\nolimits}
\def\Spec{\mathop{\mathrm{Spec}}}
\def\Hom{\mathop{\mathrm{Hom}}\nolimits}
\def\Ext{\mathop{\mathrm{Ext}}\nolimits}
\def\SheafHom{\mathop{\mathcal{H}\!\mathit{om}}\nolimits}
\def\SheafExt{\mathop{\mathcal{E}\!\mathit{xt}}\nolimits}
\def\Sch{\mathit{Sch}}
\def\Mor{\mathop{\mathrm{Mor}}\nolimits}
\def\Ob{\mathop{\mathrm{Ob}}\nolimits}
\def\Sh{\mathop{\mathit{Sh}}\nolimits}
\def\NL{\mathop{N\!L}\nolimits}
\def\CH{\mathop{\mathrm{CH}}\nolimits}
\def\proetale{{pro\text{-}\acute{e}tale}}
\def\etale{{\acute{e}tale}}
\def\QCoh{\mathit{QCoh}}
\def\Ker{\mathop{\mathrm{Ker}}}
\def\Im{\mathop{\mathrm{Im}}}
\def\Coker{\mathop{\mathrm{Coker}}}
\def\Coim{\mathop{\mathrm{Coim}}}

% Boxtimes
%
\DeclareMathSymbol{\boxtimes}{\mathbin}{AMSa}{"02}

%
% Macros for moduli stacks/spaces
%
\def\QCohstack{\mathcal{QC}\!\mathit{oh}}
\def\Cohstack{\mathcal{C}\!\mathit{oh}}
\def\Spacesstack{\mathcal{S}\!\mathit{paces}}
\def\Quotfunctor{\mathrm{Quot}}
\def\Hilbfunctor{\mathrm{Hilb}}
\def\Curvesstack{\mathcal{C}\!\mathit{urves}}
\def\Polarizedstack{\mathcal{P}\!\mathit{olarized}}
\def\Complexesstack{\mathcal{C}\!\mathit{omplexes}}
% \Pic is the operator that assigns to X its picard group, usage \Pic(X)
% \Picardstack_{X/B} denotes the Picard stack of X over B
% \Picardfunctor_{X/B} denotes the Picard functor of X over B
\def\Pic{\mathop{\mathrm{Pic}}\nolimits}
\def\Picardstack{\mathcal{P}\!\mathit{ic}}
\def\Picardfunctor{\mathrm{Pic}}
\def\Deformationcategory{\mathcal{D}\!\mathit{ef}}

\setlength{\emergencystretch}{3em}
\begin{document}
\title[Stacks Project, Tag 0F89]{389 --- A finite family of boundary ideals detects the resolution property}
\maketitle
\noindent Copyright (C) 2005--2025 Johan de Jong. Stacks Project authors. Source: \href{https://stacks.math.columbia.edu/tag/0F89}{Tag 0F89}.

\noindent Editorial difficulty note (not author text). Difficulty H2: The proof must replace the requirement to resolve every finite-type quasi-coherent sheaf by finitely many specified ideal quotients. It descends the boundary ideals and vector bundles to a Noetherian model, then extends generators across boundary ideal powers and uses tensor powers of the chosen bundles to resolve an arbitrary coherent sheaf..

\noindent Editorial provenance note (not author text). History: excerpt from revision \texttt{a0444\allowbreak 6e57e\allowbreak c1fbc\allowbreak 252a8\allowbreak 71afc\allowbreak ec775\allowbreak 2fb28\allowbreak 07b14}, perfect.tex, lines 9349--9409. Reference presentation was adapted; original-author context and core support blocks were retained or added. The mathematical wording is unchanged. Licensed under GNU FDL 1.2 or later; no invariant sections or cover texts. See ../licenses/COPYING. Context blocks reproduce author definitions and supporting results; they are not additional counted problems.

\begin{definition}
\label{definition-resolution-property}
Let $X$ be a scheme. We say $X$ has the {\it resolution property}
if every quasi-coherent $\mathcal{O}_X$-module of finite type
is the quotient of a finite locally free $\mathcal{O}_X$-module.
\end{definition}

\begin{lemma}
\label{lemma-resolution-property-finite-number}
Let $X$ be a scheme. Suppose given
\begin{enumerate}
\item a finite affine open covering $X = U_1 \cup \ldots \cup U_m$
\item finite type quasi-coherent ideals $\mathcal{I}_j$
with $V(\mathcal{I}_j) = X \setminus U_j$
\end{enumerate}
Then $X$ has the resolution property if and only if $\mathcal{I}_j$
is the quotient of a finite locally free $\mathcal{O}_X$-module
for $j = 1, \ldots, m$.
\end{lemma}

\begin{proof}
One direction of the lemma is trivial. For the other, 
say $\mathcal{E}_j \to \mathcal{I}_j$ is a surjection with
$\mathcal{E}_j$ finite locally free. In the next paragraph, we
reduce to the Noetherian case; we suggest the reader skip it.

\medskip\noindent
The first observation is that $U_j \cap U_{j'}$ is quasi-compact
as the complement of the zero scheme of the quasi-coherent finite type ideal
$\mathcal{I}_{j'}|{U_j}$ on the affine scheme $U_j$, see
Properties, Lemma \href{https://stacks.math.columbia.edu/tag/01PH}{lemma quasi coherent finite type ideals (Tag 01PH)}.
Hence $X$ is quasi-compact and quasi-separated, see
Schemes, Lemma \href{https://stacks.math.columbia.edu/tag/01KO}{lemma characterize quasi separated (Tag 01KO)}.
By Limits, Proposition \href{https://stacks.math.columbia.edu/tag/01ZA}{proposition approximate (Tag 01ZA)}
we can write $X = \lim X_i$ as the limit of a direct system of
Noetherian schemes with affine transition morphisms. For each $j$
we can find an $i$ and a finite locally free $\mathcal{O}_{X_i}$-module
$\mathcal{E}_{i, j}$ pulling back to $\mathcal{E}_j$, see
Limits, Lemma \href{https://stacks.math.columbia.edu/tag/0B8W}{lemma descend invertible modules (Tag 0B8W)}.
After increasing $i$ we may assume that the composition
$\mathcal{E}_j \to \mathcal{I}_j \to \mathcal{O}_X$ is the
pullback of a map $\mathcal{E}_{i, j} \to \mathcal{O}_{X_i}$, see
Limits, Lemma \href{https://stacks.math.columbia.edu/tag/01ZR}{lemma descend modules finite presentation (Tag 01ZR)}.
Denote $\mathcal{I}_{i, j} \subset \mathcal{O}_{X_i}$ the image
of this map; this is a quasi-coherent ideal sheaf on the Noetherian
scheme $X_i$ whose pullback to $X$ is $\mathcal{I}_j$.
Denoting $U_{i, j} \subset X_i$ the complementary opens, we
may assume these are affine for all $i, j$, see
Limits, Lemma \href{https://stacks.math.columbia.edu/tag/01Z6}{lemma limit affine (Tag 01Z6)}.
If we can prove the lemma for the opens $U_{i, j}$ and the
ideal sheaves $\mathcal{I}_{i, j}$ on $X_i$ then $X$, being affine
over $X_i$, will have the resolution property by
Lemma \href{https://stacks.math.columbia.edu/tag/0F88}{lemma resolution property goes up affine (Tag 0F88)}.
In this way we reduce to the case of a Noetherian scheme.

\medskip\noindent
Assume $X$ is Noetherian. For every coherent module $\mathcal{F}$
we can choose a finite list of sections $s_{jk} \in \mathcal{F}(U_j)$,
$k = 1, \ldots, e_j$ which generate the restriction of $\mathcal{F}$ to $U_j$.
By Cohomology of Schemes, Lemma \href{https://stacks.math.columbia.edu/tag/01YB}{lemma homs over open (Tag 01YB)}
we can extend $s_{jk}$ to a map
$s'_{jk} : \mathcal{I}_i^{n_{jk}} \to \mathcal{F}$ for some $n_{jk} \geq 1$.
Then we can consider the compositions
$$
\mathcal{E}_j^{\otimes n_{jk}} \to \mathcal{I}_j^{n_{jk}} \to \mathcal{F}
$$
to conclude.
\end{proof}
\end{document}
